% Relatorio do trabalho final de LFA - compilar com XeLaTeX ou Tectonic:
%   tectonic relatorio.tex      (ou: xelatex relatorio.tex, duas vezes)
\documentclass[12pt,a4paper]{article}

\usepackage{fontspec}
\setmainfont{texgyretermes}[
  Extension = .otf, UprightFont = *-regular, BoldFont = *-bold,
  ItalicFont = *-italic, BoldItalicFont = *-bolditalic]
\usepackage{unicode-math}
\setmathfont{texgyretermes-math.otf}
\usepackage[brazilian]{babel}
\usepackage[a4paper,top=3cm,left=3cm,right=2cm,bottom=2cm]{geometry}
\usepackage{setspace}
\usepackage{indentfirst}
\usepackage{titlesec}
\usepackage{fancyhdr}
\usepackage{booktabs}
\usepackage{array}
\usepackage{enumitem}
\usepackage{xcolor}
\usepackage{listings}
\usepackage{caption}
\usepackage{float}
\usepackage{pgfplots}
\pgfplotsset{compat=1.18}
\usepackage[edges]{forest}
\usepackage[hidelinks]{hyperref}

\onehalfspacing
\setlength{\emergencystretch}{3em}

% figuras e tabelas: mais espaco por pagina e paginas de figuras alinhadas no topo
\renewcommand{\topfraction}{0.9}
\renewcommand{\bottomfraction}{0.8}
\renewcommand{\textfraction}{0.07}
\renewcommand{\floatpagefraction}{0.75}
\makeatletter
\setlength{\@fptop}{0pt}
\setlength{\@fpsep}{12pt}
\setlength{\@fpbot}{0pt plus 1fil}
\makeatother
\setlength{\parindent}{1.25cm}
\setlength{\parskip}{0pt}

% secoes no estilo ABNT: "1 INTRODUCAO"
\titleformat{\section}{\normalfont\bfseries}{\thesection}{1em}{\MakeUppercase}
\titleformat{\subsection}{\normalfont\bfseries}{\thesubsection}{1em}{\MakeUppercase}
\titlespacing*{\section}{0pt}{18pt}{12pt}
\titlespacing*{\subsection}{0pt}{12pt}{6pt}

% numero da pagina no canto superior direito
\pagestyle{fancy}
\fancyhf{}
\fancyhead[R]{\small\thepage}
\renewcommand{\headrulewidth}{0pt}

% legendas: "Figura 1 – Titulo" em cima e a fonte embaixo
\captionsetup{labelsep=endash,font={small},justification=centering,skip=6pt}
\DeclareCaptionType{quadro}[Quadro][Lista de quadros]
\captionsetup[quadro]{labelsep=endash,font={small}}
\newcommand{\fonte}[1][Elaborado pelos autores (2026).]{%
  \par\vspace{3pt}{\footnotesize Fonte: #1}\par}

\lstset{
  basicstyle=\ttfamily\footnotesize,
  columns=fullflexible,
  keepspaces=true,
  frame=single,
  rulecolor=\color{black!40},
  xleftmargin=4pt, xrightmargin=4pt,
  aboveskip=4pt, belowskip=2pt,
  showstringspaces=false,
  upquote=true
}

\newcommand{\code}[1]{\texttt{#1}}
\newcommand{\miniPascal}{miniPascal}

\begin{document}

% ----------------------------------------------------------------------
% CAPA
% ----------------------------------------------------------------------
\begin{titlepage}
\begin{center}
\singlespacing
\textbf{UNIVERSIDADE DO ESTADO DE SANTA CATARINA -- UDESC}\\
\textbf{CENTRO DE CIÊNCIAS TECNOLÓGICAS -- CCT}\\
\textbf{DEPARTAMENTO DE CIÊNCIA DA COMPUTAÇÃO}

\vspace{4cm}
\textbf{VINICIUS MOREIRA KRÜGER}\\
\textbf{GUSTAVO BADA ZANATTA}\\
\textbf{LUCAS DA SILVA}

\vspace{4cm}
\textbf{RECONHECIMENTO DA LINGUAGEM MINIPASCAL COM O ALGORITMO CYK}

\vspace{0.6cm}
Relatório do trabalho final da disciplina de Linguagens Formais e Autômatos (LFA) -- 2026/2\\[2pt]
Professor: Leandro Israel Pinto

\vfill
\textbf{JOINVILLE}\\
\textbf{2026}
\end{center}
\end{titlepage}
\setcounter{page}{2}

% ----------------------------------------------------------------------
\section{Introdução}
% ----------------------------------------------------------------------

Este relatório apresenta o trabalho final da disciplina de Linguagens Formais e Autômatos (LFA),
feito por Vinicius Moreira Krüger, Gustavo Bada Zanatta e Lucas da Silva, estudantes de Ciência
da Computação da UDESC, em Joinville, no semestre de 2026/2.

O objetivo do trabalho é desenvolver um programa capaz de reconhecer a linguagem \miniPascal{},
isto é, dizer se um arquivo de entrada pertence ou não à linguagem, mostrando os passos do
algoritmo de reconhecimento. O programa foi escrito em C e funciona em duas etapas: a análise
léxica, feita com a ferramenta \textit{flex}, transforma o texto em uma sequência de
\textit{tokens}; depois, o algoritmo de Cocke-Younger-Kasami (CYK), visto em aula, decide se
essa sequência é gerada pela Gramática Livre do Contexto (GLC) do \miniPascal{}.

Para usar o CYK, a gramática precisa estar na Forma Normal de Chomsky (FNC). Por isso o
programa também implementa os algoritmos de simplificação de gramáticas (exclusão de produções
vazias, de produções unitárias e de símbolos inúteis) e a conversão para a FNC, e mostra cada
um desses passos. Todos os testes deram o resultado esperado: os 3 programas válidos foram
aceitos e os 2 inválidos foram rejeitados.

% ----------------------------------------------------------------------
\section{A linguagem miniPascal}
% ----------------------------------------------------------------------

A linguagem foi definida a partir do trecho do enunciado, do arquivo \code{teste.pas} e do
\code{pascal.l} enviados pelo professor, e completada com as construções do Pascal que esse tipo
de programa normalmente usa. Um programa \miniPascal{} pode começar com um cabeçalho opcional
\code{program nome;}, seguido de declarações e, se houver, de um bloco principal
\code{begin}~\dots~\code{end.}. As declarações podem ser:

\begin{itemize}[nosep]
  \item variáveis globais, escritas como no enunciado (\code{myglobalvar: real;}) ou dentro de
        uma seção \code{var}, com uma ou mais variáveis separadas por vírgula (\code{a, b: integer;});
  \item uma \code{procedure}, com lista de parâmetros separados por \code{;} (cada grupo pode ter
        vários nomes e ser passado por referência com \code{var}), declarações locais opcionais
        (seção \code{var} e sub-rotinas aninhadas) e um bloco \code{begin}~\dots~\code{end;};
  \item uma \code{function}, igual à \code{procedure}, mas com o tipo de retorno depois dos
        parâmetros. O valor de retorno é atribuído a \code{Result}.
\end{itemize}

Os tipos são \code{integer}, \code{real}, \code{boolean}, \code{char} e \code{string}. Os
comandos são a atribuição (\code{:=}), a chamada de procedimento (com ou sem argumentos), o
\code{if}~\dots~\code{then} com \code{else} opcional, o \code{while}~\dots~\code{do}, o
\code{for}~\dots~\code{to}/\code{downto}~\dots~\code{do}, o \code{repeat}~\dots~\code{until} e o
bloco \code{begin}~\dots~\code{end}, que pode ser vazio. As expressões seguem a precedência do
Pascal: \code{not}; depois \code{* / div mod and}; depois \code{+ - or}; e por último os
relacionais \code{= <> < > <= >=}. Também há sinal unário, parênteses, chamadas de função,
números inteiros e reais, textos entre apóstrofos e \code{true}/\code{false}.

Como em Pascal, maiúsculas e minúsculas são iguais (\code{Result} e \code{result} são o mesmo
identificador), o \code{;} separa comandos (o último antes do \code{end} pode ou não ter
\code{;}) e há três formas de comentário: \code{\{ \}}, \code{(* *)} e \code{//}. Ficaram de
fora os tipos estruturados (\code{array}, \code{record}), \code{const}, \code{type},
\code{case} e \code{goto}.

\subsection{Por que uma gramática livre do contexto}

A linguagem não é regular. Os blocos \code{begin}~\dots~\code{end} podem ser aninhados em
qualquer profundidade e precisam estar balanceados, o mesmo duplo balanceamento de
$\{a^n b^n \mid n \ge 0\}$ visto em aula, que nenhum autômato finito reconhece. Já a forma de
cada \textit{token} (identificador, número, palavra reservada) é regular, e por isso fica a cargo
do \textit{flex}, que gera um autômato finito determinístico.

Algumas regras de um compilador de verdade, como ``toda variável usada foi declarada'', não são
livres do contexto: elas são parecidas com a linguagem $\{ww\}$, que só uma gramática sensível ao
contexto gera. Essas regras são semânticas e ficam fora do reconhecimento. Por exemplo, o
\code{teste.pas} usa a variável \code{var1} sem declará-la e mesmo assim está sintaticamente
correto.

\subsection{Gramática}

O Quadro~\ref{q:glc} mostra a GLC $G = (V, T, P, \mathit{Programa})$ usada, do jeito que
está no arquivo \code{gramatica.txt}. Variáveis começam com letra maiúscula, os terminais são as
palavras reservadas, os símbolos e os tipos de \textit{token} \code{id}, \code{num\_int},
\code{num\_real} e \code{texto}, e \code{\&} é a palavra vazia. Ela tem 37 variáveis, 49
terminais e 92 produções. As produções de \code{ExprSimples}, \code{Termo} e \code{Fator} dão a
precedência dos operadores.

\begin{quadro}[htbp]
\caption{GLC do miniPascal (arquivo \code{gramatica.txt})}\label{q:glc}
\begin{lstlisting}
Programa       -> Cabecalho Corpo | Corpo
Cabecalho      -> program id ;
Corpo          -> ListaDecl | ListaDecl BlocoPrincipal | BlocoPrincipal
BlocoPrincipal -> begin ListaCmd end .
ListaDecl      -> Decl | Decl ListaDecl
Decl           -> DeclVar | SecaoVar | DeclProc | DeclFunc
DeclVar        -> ListaId : Tipo ;
ListaId        -> id | id , ListaId
Tipo           -> integer | real | boolean | char | string
SecaoVar       -> var ListaVar
ListaVar       -> DeclVar | DeclVar ListaVar
DeclProc       -> procedure id Parametros ; DeclLocais Bloco ;
DeclFunc       -> function id Parametros : Tipo ; DeclLocais Bloco ;
Parametros     -> ( ListaParam ) | &
ListaParam     -> Param | Param ; ListaParam
Param          -> ListaId : Tipo | var ListaId : Tipo
DeclLocais     -> DeclLocal DeclLocais | &
DeclLocal      -> SecaoVar | DeclProc | DeclFunc
Bloco          -> begin ListaCmd end
ListaCmd       -> Cmd | Cmd ; ListaCmd
Cmd            -> Atribuicao | Chamada | CmdIf | CmdWhile | CmdFor | CmdRepeat
                  | Bloco | &
Atribuicao     -> id := Expr
Chamada        -> id | id ( ListaExpr )
CmdIf          -> if Expr then Cmd | if Expr then Cmd else Cmd
CmdWhile       -> while Expr do Cmd
CmdFor         -> for id := Expr Direcao Expr do Cmd
Direcao        -> to | downto
CmdRepeat      -> repeat ListaCmd until Expr
Expr           -> ExprSimples | ExprSimples OpRel ExprSimples
ExprSimples    -> Termo | Sinal Termo | ExprSimples OpAd Termo
Termo          -> Fator | Termo OpMul Fator
Fator          -> id | id ( ListaExpr ) | num_int | num_real | texto | true | false
                  | ( Expr ) | not Fator
ListaExpr      -> Expr | Expr , ListaExpr
Sinal          -> + | -
OpRel          -> = | <> | < | > | <= | >=
OpAd           -> + | - | or
OpMul          -> * | / | div | mod | and
\end{lstlisting}
\fonte
\end{quadro}

\subsection{Ambiguidade}

Esta gramática é ambígua, como a gramática de expressões vista em aula. O caso clássico é o
``\textit{else} pendurado'': em \code{if a > 0 then if b > 0 then x := 1 else x := 2}, o
\code{else} pode pertencer a qualquer um dos dois \code{if}, e cada leitura tem a sua árvore de
derivação. Também há ambiguidade nas variáveis globais: em \code{var a: integer; x: real;}, a
declaração de \code{x} pode fazer parte da seção \code{var} ou ser uma declaração separada, o que
não muda o significado.

Seria possível escrever uma gramática não ambígua para a mesma linguagem, separando os comandos
\code{if} com e sem \code{else}, mas ela ficaria maior e mais difícil de ler. Como o CYK funciona
com qualquer GLC, inclusive ambígua, mantivemos a versão mais simples. Ele só verifica se
\code{Programa} pertence a $V_{1,n}$; quando existe mais de uma árvore, o programa mostra uma
delas (Seção~\ref{s:ambiguidade}).

% ----------------------------------------------------------------------
\section{Escolha do algoritmo}
% ----------------------------------------------------------------------

Foram vistos em aula três reconhecedores para linguagens livres do contexto. A
Tabela~\ref{t:algoritmos} resume as diferenças entre eles.

\begin{table}[H]
\caption{Comparação dos algoritmos de reconhecimento vistos em aula}\label{t:algoritmos}
\centering\small
\begin{tabular}{@{}>{\raggedright}p{2.6cm}>{\raggedright}p{3.6cm}>{\raggedright}p{2.6cm}>{\raggedright\arraybackslash}p{4.8cm}@{}}
\toprule
Algoritmo & Gramática exigida & Custo & Observação \\
\midrule
Autômato com pilha descendente & sem recursão à esquerda & $k^{|w|}$ & não determinístico; simulado com tentativa e erro \\
CYK & Forma Normal de Chomsky & $|w|^3$ & \textit{bottom-up}; tabela triangular \\
Earley & qualquer GLC & $|w|^3$ ($|w|^2$ se não ambígua) & \textit{top-down}; conjuntos $D_0 \dots D_n$ \\
\bottomrule
\end{tabular}
\fonte
\end{table}

Escolhemos o CYK pelos seguintes motivos:

\begin{itemize}[nosep]
  \item usa o conteúdo das aulas de simplificação e formas normais: a GLC escrita à mão é
        convertida automaticamente para a FNC, e o programa mostra cada etapa;
  \item o custo é sempre $O(|w|^3)$, mesmo com a gramática ambígua. O autômato com pilha
        descendente pode ter custo exponencial e exigiria tirar a recursão à esquerda de
        \code{ExprSimples} e \code{Termo};
  \item a tabela triangular $V_{r,s}$ é fácil de mostrar passo a passo, igual ao exemplo dos slides;
  \item é simples de implementar, apenas três laços aninhados.
\end{itemize}

O Earley evitaria a conversão para a FNC, mas é mais trabalhoso de implementar e de ilustrar.

% ----------------------------------------------------------------------
\section{Implementação}
% ----------------------------------------------------------------------

\subsection{Organização do código}

O trabalho foi feito em C. Na pasta de entrega, \code{src} tem o código e \code{testes} tem os
programas de teste. Os arquivos enviados pelo professor foram usados como base. Os arquivos de
\code{src} são:

\begin{itemize}[nosep]
  \item \code{pascal.l} e \code{tokens.h}: analisador léxico (\textit{flex}). O \code{lex.yy.c} é
        gerado a partir dele;
  \item \code{gramatica.txt}: a GLC do Quadro~\ref{q:glc}, lida pelo programa ao iniciar;
  \item \code{gramatica.c}: leitura da GLC, simplificação e conversão para a FNC;
  \item \code{cyk.c}: o algoritmo CYK e a árvore de derivação;
  \item \code{main.c}: lê o arquivo de entrada, chama as etapas e mostra o resultado.
\end{itemize}

A gramática fica em vetores globais: \code{esq[p]} é a variável do lado esquerdo da produção
$p$, \code{dir[p]} são os símbolos do lado direito e \code{tam[p]} é a quantidade deles
(0 para a palavra vazia). Os símbolos são guardados pelo nome em \code{simbolos[]}.

\subsection{Análise léxica}

O \code{pascal.l} do professor foi ampliado. Cada padrão é uma expressão regular, e o
\textit{flex} junta todos em um único autômato finito determinístico. O léxico devolve só sete
tipos de \textit{token}: palavra reservada, identificador, número inteiro, número real, texto
entre apóstrofos, símbolo (\code{:= ; ( ) +} etc.) e erro. Comentários e espaços são ignorados,
e a opção \code{caseless} faz o autômato ignorar maiúsculas e minúsculas.

No \code{main.c}, cada \textit{token} vira um terminal da gramática: a palavra reservada e o
símbolo viram eles mesmos (\code{begin}, \code{:=}), e os demais viram \code{id},
\code{num\_int}, \code{num\_real} ou \code{texto}. Assim, a palavra $w$ que o CYK reconhece é a
sequência de \textit{tokens}, e não de caracteres: \code{myglobalvar: real;} vira
$w = $ \code{id : real ;}, com $|w| = 4$. Um caractere que não pertence ao alfabeto, como
\code{@}, gera o \textit{token} de erro e o arquivo é rejeitado já nesta etapa.

\subsection{Simplificação e Forma Normal de Chomsky}

O programa aplica os algoritmos dos slides na ordem recomendada em ``Simplificações
Combinadas'': primeiro exclui as produções vazias, depois as produções da forma $A \to B$ e por
último os símbolos inúteis. Em seguida converte o resultado para a FNC. A Tabela~\ref{t:fnc}
mostra o tamanho da gramática depois de cada etapa.

\begin{table}[H]
\caption{Tamanho da gramática em cada etapa da conversão}\label{t:fnc}
\centering
\begin{tabular}{@{}lccc@{}}
\toprule
Etapa & Variáveis & Terminais & Produções \\
\midrule
GLC original              & 37  & 49 & 92 \\
Sem produções vazias      & 37  & 49 & 108 \\
Sem produções $A \to B$   & 37  & 49 & 228 \\
Sem símbolos inúteis      & 28  & 49 & 204 \\
Forma Normal de Chomsky   & 118 & 49 & 294 \\
\bottomrule
\end{tabular}
\fonte
\end{table}

Na exclusão das produções vazias, o programa encontra as variáveis que geram a palavra vazia,
$V_\varepsilon = \{\mathit{ListaCmd},\allowbreak \mathit{Parametros},\allowbreak
\mathit{DeclLocais},\allowbreak \mathit{Cmd}\}$, e para cada ocorrência delas cria uma versão da
produção sem ela. Uma das produções novas é a versão sem parâmetros e sem declarações locais:
\begin{center}\code{DeclProc -> procedure id ; Bloco ;}\end{center}
Como $\mathit{Programa} \notin V_\varepsilon$, a palavra vazia não pertence à linguagem e nenhuma
produção $S \to \varepsilon$ é incluída.

Na exclusão das produções unitárias, o programa calcula o fecho de cada variável, por exemplo
$\textit{FECHO-Expr} = \{\mathit{ExprSimples},\allowbreak \mathit{Termo},\allowbreak
\mathit{Fator}\}$, e copia as produções das variáveis do fecho. Depois disso, nove variáveis
(\code{SecaoVar}, \code{DeclProc}, \code{DeclFunc}, \code{Atribuicao}, \code{Chamada},
\code{CmdIf}, \code{CmdWhile}, \code{CmdFor} e \code{CmdRepeat}) só apareciam em produções
unitárias, deixam de ser atingíveis a partir de \code{Programa} e são retiradas na etapa de
símbolos inúteis. É por isso que essa etapa vem por último.

Na conversão para a FNC, cada terminal que aparece em produção de tamanho 2 ou mais ganha uma
variável própria (23 ao todo), como $T(\mathtt{;}) \to \mathtt{;}$, e as produções de tamanho 3
ou mais são quebradas em cascata, começando pelo fim:
\begin{gather*}
\mathit{DeclVar} \to \mathit{ListaId}\, T(\mathtt{:})\, \mathit{Tipo}\, T(\mathtt{;})
\\
\Downarrow
\\
C_4 \to \mathit{Tipo}\, T(\mathtt{;}),\qquad C_5 \to T(\mathtt{:})\, C_4,\qquad
\mathit{DeclVar} \to \mathit{ListaId}\, C_5
\end{gather*}
Antes de criar uma variável $C_i$, o programa procura se já existe uma com a mesma produção e,
se existir, reaproveita. Com isso a FNC ficou com 67 variáveis de cascata em vez de 370. Ao
final, todas as 294 produções têm a forma $A \to BC$ ou $A \to a$. A saída de cada etapa aparece
com a opção \code{-g}.

\subsection{Algoritmo CYK}

O CYK foi implementado como no pseudocódigo dos slides (Quadro~\ref{q:cyk}). A célula
$V_{r,s}$ guarda as variáveis que geram a subpalavra que começa no \textit{token} $r$ e tem
tamanho $s$. A palavra é aceita se o símbolo inicial \code{Programa} pertence a $V_{1,n}$.

\begin{quadro}[H]
\caption{Pseudocódigo do CYK implementado em \code{cyk.c}}\label{q:cyk}
\begin{lstlisting}
para r de 1 ate n:                      // linha s = 1: producoes A -> a
    V[r][1] = { A | A -> a_r em P }
para s de 2 ate n:                      // producoes A -> BC
    para r de 1 ate n - s + 1:
        V[r][s] = {}
        para k de 1 ate s - 1:
            V[r][s] = V[r][s] U { A | A -> BC em P,
                                      B em V[r][k] e C em V[r+k][s-k] }
aceita se Programa pertence a V[1][n]
\end{lstlisting}
\fonte
\end{quadro}

A tabela é um vetor de \code{char} com uma posição para cada célula e cada símbolo: o valor 1
indica que a variável pertence à célula. O programa também guarda se cada célula tem alguma
variável. Quando $V_{r,k}$ ou $V_{r+k,s-k}$ está vazia, nenhuma produção pode ser usada e a
combinação $(s, r, k)$ é pulada. O programa conta as combinações visitadas, quantas delas
tinham as duas células preenchidas, e mede o tempo do CYK.

Se a palavra for aceita, a árvore de derivação é reconstruída a partir da tabela. Partindo de
\code{Programa} em $V_{1,n}$, procura-se um $k$ e uma produção $A \to BC$ com $B \in V_{r,k}$ e
$C \in V_{r+k,s-k}$, e repete-se o processo nos dois filhos. Na impressão, as variáveis criadas
na FNC ($T(a)$ e $C_i$) são omitidas, e a árvore fica parecida com a da gramática original.

% ----------------------------------------------------------------------
\section{Compilação e execução}
% ----------------------------------------------------------------------

O programa é compilado pela linha de comando, dentro da pasta \code{src}. No Windows, os
arquivos \code{COMPILAR.bat} e \code{RUN FLEX.bat} executam esses comandos:

\begin{lstlisting}[frame=none]
flex -o lex.yy.c pascal.l      (so se o pascal.l for alterado)
gcc -Wall -O2 -o main main.c gramatica.c cyk.c lex.yy.c
./main [opcoes] arquivo.pas
\end{lstlisting}

O código compila sem avisos no GCC 15 (MinGW, Windows) e no GCC 14 (Linux). Sem opções, o
programa mostra a lista de \textit{tokens}, o resumo do CYK e o resultado. As opções da
Tabela~\ref{t:opcoes} mostram os outros passos. O código de saída é 0 quando o arquivo é aceito
e 1 quando é rejeitado.

\begin{table}[H]
\caption{Opções do programa}\label{t:opcoes}
\centering
\begin{tabular}{@{}cl@{}}
\toprule
Opção & O que mostra \\
\midrule
\code{-g} & cada etapa da simplificação e da conversão para a FNC \\
\code{-p} & o CYK passo a passo: cada célula, com os valores de $s$, $r$ e $k$ \\
\code{-a} & a árvore de derivação \\
\bottomrule
\end{tabular}
\fonte
\end{table}

% ----------------------------------------------------------------------
\section{Ilustração do reconhecimento}
% ----------------------------------------------------------------------

\subsection{Exemplo passo a passo}

Para mostrar o funcionamento do CYK com uma tabela que cabe na página, usamos o menor programa
possível, \code{x: integer;} (arquivo \code{src/exemplo.pas}), que vira a palavra
$w = $ \code{id : integer ;} com $n = 4$. A Figura~\ref{f:tabela} mostra a tabela gerada pelo
programa com a opção \code{-p}, no mesmo formato dos slides.

\begin{figure}[htbp]
\caption{Tabela do CYK para \code{x: integer;}}\label{f:tabela}
\centering\small
\renewcommand{\arraystretch}{1.3}
\begin{tabular}{r|>{\centering\arraybackslash}p{3.4cm}|>{\centering\arraybackslash}p{2.1cm}|>{\centering\arraybackslash}p{2.1cm}|>{\centering\arraybackslash}p{2.1cm}|}
\cline{2-2}
$s=4$ & Programa, Corpo, ListaDecl, Decl, DeclVar, ListaVar \\
\cline{2-3}
$s=3$ & ListaParam, Param, $C_{10}$ & $C_5$ \\
\cline{2-4}
$s=2$ & $\emptyset$ & $C_9$ & $C_4$ \\
\cline{2-5}
$s=1$ & ListaCmd, ListaId, Cmd, Expr, ListaExpr, ExprSimples, Termo, Fator, $T(\mathtt{id})$ & $T(\mathtt{:})$ & Tipo & ListaCmd, $T(\mathtt{;})$ \\
\cline{2-5}
\multicolumn{1}{r}{} & \multicolumn{1}{c}{\code{id}} & \multicolumn{1}{c}{\code{:}} & \multicolumn{1}{c}{\code{integer}} & \multicolumn{1}{c}{\code{;}} \\
\multicolumn{1}{r}{} & \multicolumn{1}{c}{$r=1$} & \multicolumn{1}{c}{$r=2$} & \multicolumn{1}{c}{$r=3$} & \multicolumn{1}{c}{$r=4$} \\
\end{tabular}
\fonte
\end{figure}

A linha $s=1$ recebe as variáveis com produção $A \to a_r$. O \code{id} aparece em várias
porque, depois da exclusão das produções unitárias, muitas variáveis passaram a ter a produção
\code{-> id}: um identificador sozinho pode ser uma expressão (\code{Expr}, \code{Termo},
\code{Fator}), uma lista de nomes (\code{ListaId}) ou uma chamada de procedimento sem argumentos
(\code{Cmd}, \code{ListaCmd}). As demais linhas vêm das produções $A \to BC$. O
Quadro~\ref{q:passos} mostra a saída da opção \code{-p} para as linhas 3 e 4.

\begin{quadro}[htbp]
\caption{Trecho da saída da opção \code{-p} para \code{x: integer;}}\label{q:passos}
\begin{lstlisting}[basicstyle=\ttfamily\scriptsize]
s=3 r=1 k=1: B em V[1][1] = {ListaCmd, ListaId, Cmd, Expr, ListaExpr, ExprSimples, Termo,
                            Fator, T(id)}, C em V[2][2] = {C9}
    Param -> ListaId C9
    C10 -> ListaId C9
    ListaParam -> ListaId C9
s=3 r=1 k=2: B em V[1][2] = {}, C em V[3][1] = {Tipo}
V[1][3] = {ListaParam, Param, C10}

s=3 r=2 k=1: B em V[2][1] = {T(:)}, C em V[3][2] = {C4}
    C5 -> T(:) C4
s=3 r=2 k=2: B em V[2][2] = {C9}, C em V[4][1] = {ListaCmd, T(;)}
V[2][3] = {C5}

s=4 r=1 k=1: B em V[1][1] = {ListaCmd, ListaId, Cmd, Expr, ListaExpr, ExprSimples, Termo,
                            Fator, T(id)}, C em V[2][3] = {C5}
    DeclVar -> ListaId C5
    Programa -> ListaId C5
    Corpo -> ListaId C5
    ListaDecl -> ListaId C5
    Decl -> ListaId C5
    ListaVar -> ListaId C5
s=4 r=1 k=2: B em V[1][2] = {}, C em V[3][2] = {C4}
s=4 r=1 k=3: B em V[1][3] = {ListaParam, Param, C10}, C em V[4][1] = {ListaCmd, T(;)}
V[1][4] = {Programa, Corpo, ListaDecl, Decl, DeclVar, ListaVar}

4 tokens, 10 combinacoes (s, r, k), 8 com as duas celulas preenchidas, 0.0 ms
ACEITA: Programa pertence a V[1][4]
\end{lstlisting}
\fonte
\end{quadro}

Como \code{Programa} está em $V_{1,4}$, a palavra é aceita. A Figura~\ref{f:arvores} mostra a
árvore de derivação na FNC, reconstruída a partir da tabela, e a mesma árvore como o programa
a imprime, sem as variáveis auxiliares.

\begin{figure}[htbp]
\caption{Árvore de derivação de \code{x: integer;} na FNC e sem as variáveis auxiliares}\label{f:arvores}
\centering\small
\begin{forest}
for tree={font=\small, l sep=8pt, s sep=6pt, inner sep=1pt}
[Programa
  [ListaId [\code{id}]]
  [{$C_5$}
    [{$T(\mathtt{:})$} [\code{:}]]
    [{$C_4$}
      [Tipo [\code{integer}]]
      [{$T(\mathtt{;})$} [\code{;}]]]]]
\end{forest}
\hspace{1.8cm}
\begin{forest}
for tree={font=\small, l sep=8pt, s sep=10pt, inner sep=1pt}
[Programa
  [ListaId [\code{id}]]
  [\code{:}]
  [Tipo [\code{integer}]]
  [\code{;}]]
\end{forest}
\fonte
\end{figure}

\subsection{Programas do enunciado e do professor}

O trecho do enunciado (81 \textit{tokens}) e o \code{teste.pas} do professor (72
\textit{tokens}) foram aceitos. No \code{teste.pas}, o CYK visita 62.196 combinações
$(s, r, k)$, mas só 481 delas (0,8\%) têm as duas células preenchidas, porque a maioria das
subpalavras não é gerada por nenhuma variável. O Quadro~\ref{q:arvore} mostra a parte da árvore
de derivação, impressa com a opção \code{-a}, que corresponde à função \code{soma} do
\code{teste.pas}.

\begin{quadro}[htbp]
\caption{Árvore de derivação da função \code{soma} do \code{teste.pas} (opção \code{-a})}\label{q:arvore}
\begin{lstlisting}
  ListaDecl
    function
    id  soma
    Parametros
      (
      ListaParam
        Param
          ListaId
            id  a
          :
          Tipo
            integer
        ;
        ListaParam
          ListaId
            id  b
          :
          Tipo
            integer
      )
    :
    Tipo
      integer
    ;
    Bloco
      begin
      ListaCmd
        Cmd
          id  Result
          :=
          Expr
            ExprSimples
              id  a
            OpAd
              +
            Termo
              id  b
        ;
      end
    ;
\end{lstlisting}
\fonte
\end{quadro}

\subsection{Entrada ambígua}\label{s:ambiguidade}

O terceiro teste válido (\code{03\_programa\_completo.pas}) tem o comando
\code{if k mod 2 = 0 then if k > 10 then total := total - 1 else total := total + 1}, que é um
caso de \textit{else} pendurado. O programa aceita a entrada e, na árvore impressa, o
\code{else} fica ligado ao \code{if} mais próximo, que é a regra do Pascal. Isso acontece porque
a reconstrução da árvore pega a primeira divisão $k$ que funciona; o CYK em si não escolhe entre
as leituras, ele só decide se a palavra pertence à linguagem.

% ----------------------------------------------------------------------
\section{Testes e resultados}\label{s:testes}
% ----------------------------------------------------------------------

Os testes ficam em \code{testes/validos} e \code{testes/invalidos} e são executados pelo
\code{rodar\_testes.bat}, que confere se cada programa válido é aceito e cada inválido é
rejeitado, pelo código de saída do programa. A Tabela~\ref{t:testes} mostra os resultados.

\begin{table}[htbp]
\caption{Testes}\label{t:testes}
\centering\small
\begin{tabular}{@{}l>{\raggedright\arraybackslash}p{7.2cm}rl@{}}
\toprule
Arquivo & O que testa & $n$ & Resultado \\
\midrule
\code{validos/01} & \code{teste.pas} enviado pelo professor & 72 & aceito \\
\code{validos/02} & trecho do enunciado: variável global, \code{procedure} e \code{function} & 81 & aceito \\
\code{validos/03} & programa completo: \code{program}, \code{var}, listas com vírgula, parâmetro
\code{var}, sub-rotina aninhada, \code{if}/\code{else} (com \textit{else} pendurado), \code{for},
\code{repeat}, chamadas, operadores lógicos, textos e bloco principal \code{end.} & 232 & aceito \\
\code{invalidos/01} & \code{begin} sem o \code{end} correspondente & 29 & rejeitado (sintaxe) \\
\code{invalidos/02} & caractere \code{@}, que não pertence ao alfabeto & -- & rejeitado (léxico) \\
\bottomrule
\end{tabular}
\fonte
\end{table}

Os 5 testes deram o resultado esperado. O \code{invalidos/01} mostra o motivo de a linguagem
precisar de uma gramática livre do contexto: o erro só aparece porque o número de \code{begin} e
de \code{end} não bate. O \code{invalidos/02} é rejeitado antes do CYK, na análise léxica, com a
mensagem \code{Erro lexico na linha 3: '@' nao pertence a linguagem}.

\subsection{Desempenho}

Para medir o custo do CYK, geramos programas com o \code{teste.pas} repetido de 1 a 14 vezes,
de 72 até 1.008 \textit{tokens}. A Figura~\ref{f:tempo} mostra o tempo do CYK e a curva
$a \cdot n^3$ ajustada ao último ponto, e a Tabela~\ref{t:desempenho} mostra alguns valores. As
medições foram feitas no GCC 14 com \code{-O2}.

\begin{figure}[htbp]
\caption{Tempo do CYK em função do número de \textit{tokens}}\label{f:tempo}
\centering
\begin{tikzpicture}
\begin{axis}[
  width=12cm, height=6.5cm,
  xlabel={Número de \textit{tokens} ($n$)}, ylabel={Tempo (ms)},
  xmin=0, xmax=1050, ymin=0,
  grid=major, grid style={gray!25},
  legend pos=north west, legend cell align=left,
  /pgf/number format/use comma, /pgf/number format/1000 sep={.},
  tick label style={font=\small}, label style={font=\small}, legend style={font=\small}]
\addplot[only marks, mark=*, mark size=1.6pt, blue!70!black]
  table[x=n, y=tempo_ms, col sep=comma] {dados/desempenho.csv};
\addlegendentry{tempo medido}
\addplot[red!70!black, thick, dashed, domain=0:1008, samples=60] {2.188e-7*x^3};
\addlegendentry{$a \cdot n^3$}
\end{axis}
\end{tikzpicture}
\fonte
\end{figure}

\begin{table}[htbp]
\caption{Custo do CYK para entradas de tamanho crescente}\label{t:desempenho}
\centering\small
\begin{tabular}{@{}rrrr@{}}
\toprule
$n$ & Combinações $(s,r,k)$ & Com as duas células preenchidas & Tempo (ms) \\
\midrule
72    & 62.196      & 481    & 0,4 \\
288   & 3.981.264   & 2.193  & 5,4 \\
504   & 21.337.260  & 4.445  & 27,9 \\
720   & 62.207.880  & 7.453  & 61,6 \\
1.008 & 170.698.584 & 13.013 & 224,1 \\
\bottomrule
\end{tabular}
\fonte
\end{table}

% ----------------------------------------------------------------------
\section{Discussão dos resultados}
% ----------------------------------------------------------------------

O programa reconheceu corretamente todos os testes, incluindo o \code{teste.pas} e o exemplo do
enunciado, e rejeitou os dois programas com erro, um na sintaxe e outro no léxico. Isso indica
que a simplificação e a conversão para a FNC preservaram a linguagem, como garantem os teoremas
vistos em aula.

O custo seguiu o esperado. O número de combinações $(s, r, k)$ é exatamente $(n^3 - n)/6$:
para $n = 72$ são 62.196. Quando $n$ dobra de 504 para 1.008, as combinações são multiplicadas
por 8,0 e o tempo também por 8,0, que é o $2^3$ de um algoritmo $O(n^3)$, e a curva $a \cdot n^3$
da Figura~\ref{f:tempo} acompanha os pontos medidos. Já as combinações com as duas células
preenchidas foram multiplicadas só por 2,9, porque a tabela é quase toda vazia. Como o programa
pula as combinações com uma célula vazia, o tamanho da gramática pesa pouco, e o que domina o
custo são os três laços.

Para os programas dos testes o CYK leva poucos milissegundos (4,7 ms no maior, com 232
\textit{tokens}). Porém, como o custo é cúbico, um programa com 10.000 \textit{tokens} levaria
em torno de 1.000 vezes o tempo de um com 1.000, ou seja, alguns minutos. É por isso que
compiladores reais usam linguagens livres do contexto determinísticas, reconhecidas por um
autômato com pilha determinístico em tempo linear. Para isso a gramática precisaria ser
reescrita sem ambiguidade (o \textit{else} pendurado, por exemplo), enquanto o CYK funciona com
qualquer GLC, sem adaptar a gramática à mão.

A conversão para a FNC mais que triplicou a gramática, de 37 para 118 variáveis e de 92 para 294
produções. A maior parte vem das cascatas, e reaproveitar as variáveis $C_i$ reduziu as
variáveis de cascata de 370 para 67. Uma consequência da FNC é que a árvore de derivação
reconstruída tem variáveis que o programador não escreveu. Por isso a impressão omite $T(a)$ e
$C_i$, e a árvore fica parecida com a da gramática original.

Uma limitação do CYK é que ele só responde se a palavra pertence ou não à linguagem, sem
apontar o \textit{token} onde está o erro. No \code{invalidos/01}, por exemplo, o programa diz
apenas que \code{Programa} não pertence a $V_{1,29}$. Um reconhecedor \textit{top-down}, como o
Earley, saberia o ponto exato em que a entrada deixou de poder ser continuada.

% ----------------------------------------------------------------------
\section{Considerações finais}
% ----------------------------------------------------------------------

O trabalho implementou em C um reconhecedor para a linguagem \miniPascal{}, juntando duas classes
de linguagens vistas na disciplina: a análise léxica é regular e foi feita com um autômato finito
gerado pelo \textit{flex}, e a sintaxe é livre do contexto e foi reconhecida pelo algoritmo CYK.
No caminho, foram aplicados os algoritmos de simplificação de gramáticas e a conversão para a
Forma Normal de Chomsky, todos mostrados passo a passo pelo programa.

Todos os testes tiveram o resultado esperado, e o CYK lidou sem problemas com a ambiguidade da
gramática. As medições confirmaram o custo $O(n^3)$, rápido para os programas do trabalho, mas
que mostra por que compiladores preferem gramáticas determinísticas para entradas grandes.

% ----------------------------------------------------------------------
\section*{REFERÊNCIAS}
% ----------------------------------------------------------------------
\begin{flushleft}
\singlespacing
\setlength{\parindent}{0pt}
\setlength{\parskip}{10pt}

EARLEY, J. An efficient context-free parsing algorithm. \textbf{Communications of the ACM},
v. 13, n. 2, p. 94--102, 1970.

HOPCROFT, J. E.; ULLMAN, J. D.; MOTWANI, R. \textbf{Introdução à teoria de autômatos,
linguagens e computação}. Rio de Janeiro: Campus, 2002.

MENEZES, P. F. B. \textbf{Linguagens formais e autômatos}. 4. ed. Porto Alegre: Sagra
Luzzatto, 2002. (Série Livros Didáticos, n. 3).

PAXSON, V. \textbf{Flex: a fast scanner generator}. Manual do flex, versão 2.6, 2017.
\end{flushleft}

\end{document}
